\section{Configuração do firewall}

\subsection{IpTables}

Iptables é uma ferramenta de firewall para sistemas Linux que atua como uma espécie de 
“guarda de fronteira” entre o seu computador e a internet. Ela permite controlar o tráfego de rede, 
decidindo quais conexões são permitidas e quais são bloqueadas.


No iptables, existem quatro tabelas principais, 
cada uma com suas cadeias (chains) e propósitos específicos:

\subsubitem{Tabela Filter}
Essa é a tabela padrão do iptables e é usada para filtrar o tráfego de pacotes com base em regras definidas. 
As três cadeias padrão nesta tabela são:

\begin{itemize}
    \item \textbf{INPUT: } Essa cadeia lida com pacotes destinados à própria máquina (tráfego de entrada).
    \item \textbf{OUTPUT: } Essa cadeia lida com pacotes originados na própria máquina e enviados a outros 
    destinos (tráfego de saída).
    \item \textbf{FORWARD: }  Essa cadeia lida com pacotes que estão apenas passando pelo sistema, 
    atuando como um roteador.
    
\end{itemize}


\subsubitem{Tabela NAT}

Essa tabela é usada para realizar a tradução de endereços de rede, geralmente usada para redirecionamento 
de portas (port forwarding) ou para mascaramento de endereços IP (SNAT – Source Network Address Translation).
As cadeias padrão nesta tabela são:


\begin{itemize}
    \item \textbf{PREROUTING: } Essa cadeia é utilizada para modificar pacotes antes de serem roteados, 
    permitindo o redirecionamento de portas.
    \item \textbf{OUTPUT: }  Esta cadeia é usada para modificar pacotes gerados localmente antes de 
    saírem do sistema.
    \item \textbf{POSTROUTING:  }  Esta cadeia é usada para modificar pacotes gerados localmente 
    antes de saírem do sistema.
    
\end{itemize}

\subsubitem{Tabela Mangle}

A tabela mangle é usada para modificar os cabeçalhos dos pacotes. É usada principalmente para marcação de
pacotes para posterior manipulação por outras regras ou ferramentas. As cadeias padrão nesta tabela são:

\begin{itemize}
    \item \textbf{PREROUTING: } Usada para modificar pacotes antes de serem roteados.
    \item  \textbf{INPUT: } Usada para modificar pacotes de entrada antes de serem entregues localmente.
    \item  \textbf{FORWARD: } Usada para modificar pacotes que estão apenas passando pelo sistema.
    \item \textbf{OUTPUT: }  Usada para modificar pacotes de saída antes de deixarem o sistema.

    \item \textbf{POSTROUTING:  }  Usada para modificar pacotes após o roteamento
    
\end{itemize}